\item Unos astrónomos observan la cromosfera del Sol con un filtro espectral de la línea roja de hidrógeno $\text{H}_\alpha$ ($\lambda = 656.3\text{ nm}$). El filtro consiste en un dieléctrico transparente ($n = 1.378$) colocado entre dos placas de vidrio parcialmente aluminizadas.
\begin{enumerate}
    \item Encuentre el grosor mínimo $d$ del dieléctrico que permita una transmisión máxima normal de la línea $\text{H}_\alpha$.
    \item \textbf{¿Qué pasaría si?} Si aumenta la temperatura del filtro dilatando su espesor, ¿qué ocurre con la longitud de onda transmitida?
    \item ¿Qué otra longitud de onda casi visible pasará también el filtro?
\end{enumerate}